%% ma: L11-B15-0094
%% lop: 11
%% chuong: 5
%% bai: 15
%% dang: 11.15.5
%% loai: TLN
%% muc: VD
%% dap_an: "6"
%% nguon: "tai-lieu-hoc-tap-mon-toan-lop-11-theo-chuong-trinh-giao-duc-pho-thong-2018.pdf (Lớp Toán Thầy Duy, Phạm Lê Duy, THPT Chu Văn An), Chương 3 Bài 1, Phần 3, Câu 16 (tr.346)"
%% kien_thuc: "Lớp 11 Bài 15 (giới hạn của dãy số)"
%% trang_thai: chinh-thuc
%% ghi_chu: "Đã duyệt 03/11/2026."
\begin{ex}
Giới hạn dãy số $u_n=\dfrac{\sqrt{n^2+2n}-\sqrt{n^2+n}}n$ có dạng $\lim\dfrac{\dfrac an}{\sqrt{1+\dfrac bn}+\sqrt{1+\dfrac cn}}$ với $a;b;c$ là các số tự nhiên. Tính giá trị $a^2+b^2+c^2$.
\shortans{6}
\loigiai{\sloppy $\sqrt{n^2+2n}-\sqrt{n^2+n}=\dfrac n{\sqrt{n^2+2n}+\sqrt{n^2+n}}$, nên \[u_n=\dfrac1{\sqrt{n^2+2n}+\sqrt{n^2+n}}=\dfrac{\frac1n}{\sqrt{1+\frac2n}+\sqrt{1+\frac1n}}.\] Vậy $a=1$, $b=2$, $c=1$ và $a^2+b^2+c^2=6$.}
\end{ex}
